\documentclass[10pt,a4paper]{amsart}
\usepackage[margin=19mm]{geometry}
\usepackage{amsmath,amssymb,mathtools}
\usepackage[scr=rsfs,cal=euler]{mathalfa}
\usepackage{xfrac}
\usepackage[unicode,hidelinks,bookmarks=false]{hyperref}
\newcommand{\ie}{i.e.\ }
\newcommand{\cf}{cf.\ }
\newcommand{\resp}{respectively\ }
\newcommand{\Subsection}{\subsection}
\providecommand{\nobreakdash}{\nobreak\hbox{-}}
\setlength{\emergencystretch}{2em}
\multlinegap0pt
\usepackage{bm}
\usepackage{bbm}
\usepackage{dsfont}
\usepackage{xspace}
\usepackage{cancel}
\usepackage{mathtools}
\usepackage{amsthm}
\makeatletter
\@ifundefined{corollary}{\newtheorem{corollary}{Corollary}}{}
\@ifundefined{theorem}{\newtheorem{theorem}{Theorem}}{}
\makeatother
\title[Finite Volume Einstein Finsler Warped Product Manifolds of N]{Finite Volume Einstein Finsler Warped Product Manifolds of Non-positive or Non-negative Scalar Curvature}
\author{Mohammad Aqib, Hemangi Madhusudan Shah, Pankaj Kumar, Anjali Shriwastawa}
\date{Source: arXiv:2602.00601v1 (CC BY 4.0)}
\begin{document}
\maketitle

\noindent\emph{Source and licence.} Finite Volume Einstein Finsler Warped Product Manifolds of Non-positive or Non-negative Scalar Curvature. Version arXiv:2602.00601v1 (CC BY 4.0), distributed under the Creative Commons Attribution 4.0 International licence, as recorded in the arXiv licence metadata for this exact version and shown on the abstract page https://arxiv.org/abs/2602.00601v1. Full rights notice: licenses/arxiv-2602.00601-CC-BY-4.0.txt.\par\smallskip

\noindent\textit{Editorial scope.} Counted unit: the Theorem whose source label is \texttt{thm4.2} in the article (source0.tex lines 1383--1409, source label \texttt{thm4.2}), delivered together with the article's own complete original proof of it (source0.tex lines 1411--1500, one \texttt{proof} environment of 90 lines). No mathematics is written, repaired or re-derived: statement and proof are the article's own text line for line. Required context retained uncounted: f10 (lines 627--631); 1.4 (lines 658--674); 1.5 (lines 667--669); 1.6 (lines 670--672). Every definition, display and label of those lines is retained unchanged. Omitted from this package and not redistributed: 1--626, 632--657, 673--1382, 1410--1410, 1501--1707. Nothing inside a retained excerpt is omitted or altered: every retained body is delivered exactly as the article's lines are written, so no omission receipt of any kind accompanies this package. Citations: the retained excerpts carry no citation command at all, so no bibliography entry of the article is transcribed and no reference is redistributed. Reference adaptation: none was needed and none was made; every reference inside the delivered unit resolves inside it, so no unresolved reference or citation is produced. Printed slips kept literal: no printed word, symbol, hypothesis or number of the counted statement or of its proof was altered. Layout and class: the article's own class is \texttt{sn-jnl} with packages including graphicx, multirow, amsmath, amssymb, amsfonts, breqn, amsthm, mathrsfs, appendix, xcolor, textcomp, manyfoot, booktabs, algorithm; the wrapper is the lane's pinned \texttt{amsart} editorial wrapper at 10pt on A4, which restates the theorem-like environments the retained text needs and adds the article's own macro definitions that the retained text needs (none), with extra wrapper packages (bm, bbm, dsfont, xspace, cancel, mathtools). The delivered numbering is the unit's own. Assets: no retained span contains a figure environment, a graphics-inclusion command, a PDF-page inclusion, a table or a TikZ picture; no image file is copied into this package, the package declares no asset dependency and no image file is redistributed with it. Licence: version arXiv:2602.00601v1 of this article is distributed under the Creative Commons Attribution 4.0 International licence, as recorded in the arXiv licence metadata for this exact version and shown on the abstract page https://arxiv.org/abs/2602.00601v1. A full rights notice is delivered at licenses/arxiv-2602.00601-CC-BY-4.0.txt. Characters: every retained excerpt is pure ASCII, so the delivered engine, XeLaTeX with its default Latin Modern text font, reports no missing character.
		\begin{align}
		(\operatorname{Ric}^{M_1})_{ij}=\lambda g_1{_{ij}}+\frac{{n_2}}{f}H_{ij}^f,\label{f8}&&\\
			 (M_2,g_{M_2})\,\text{is  Einstein with }\, \operatorname{Ric}_{\alpha\beta}^{M_2}=\mu g_{\alpha\beta},\label{f9}&&\\
		 \mu=\lambda f^2 - f \Delta f + (n_2-1)  g_1(\nabla f,\nabla f).\label{f10}
		\end{align}

\begin{corollary}
If $M_{1}^{n_1}\times_{f}M_{2}^{n_2}$ is an Einstein Finslerian warped product manifold with 
$\operatorname{Ric}_{ab}=\lambda g_{ab}$, 
where $(M_{1}^{n_1},g_1)$ is a Riemannian manifold and $(M_2^{n_2},F_2)$ is a weakly Berwald Finsler manifold.
Then the scalar curvatures 
 of warped product, base and fiber satisfy:
\begin{equation}\label{1.4}
			\operatorname{Scal}_{M_1}=n_1\lambda- n_2 \frac{\Delta f}{f},
		\end{equation}
 \begin{equation}\label{1.5}
			\operatorname{Scal}_{M_2}=n_2\mu,
		\end{equation}
  \begin{equation}\label{1.6}
\operatorname{Scal}_{M_1}+\dfrac{\operatorname{Scal}_{M_2}}{f^2}=n_2(n_2 -1)\dfrac{\|\nabla f\|^2}{f^2}-2n_2\frac{\Delta f}{f}+\operatorname{Scal}_M .		
	\end{equation}
    
\end{corollary}

\begin{theorem}\label{thm4.2}
	Let $M=M_1^{n_1}\times_{f} M_2^{n_2}$ be an Einstein warped product Finsler manifold with 
 positive scalar curvature. 
 Let the base be a Riemannian manifold and fiber be  a weakly Berwald Finsler manifold with 
positive scalar curvature.
Suppose that 
  any of  the following conditions holds, then the 
  warping function is either subharmonic or superharmonic. Consequently, if the base manifold is compact 
  or if either $f$ attains its maximum or minimum, then
  it is a constant function. In this case,
   the warped product $M=M_1\times_f M_2$ is a weakly Berwald Finsler manifold, which is  simply a Riemannian product of two Finslerian manifolds.
	\begin{itemize}
		\item[\emph{(a)}] There exist a point $\Tilde{x_2}\in M_2$ such that 
  $\|\nabla f(x_1)\|\geq \sqrt{\frac{\operatorname{Scal}_{M_2}(\Tilde{x_2})}{n_2(n_2-1)}}$ for every $x_1$,
		
		\item[\emph{(b)}] $\operatorname{Scal}_M\leq \operatorname{Scal}_{M_1}$,
		
		\item[\emph{(c)}] $\operatorname{Scal}_M-\operatorname{Scal}_{M_2}\geq \operatorname{Scal}_{M_1}$ and $\frac{\operatorname{Scal}_{M_2}(x_2)}{\operatorname{Scal}_{M}(x_1,x_2)}\geq\frac{n_2}{(n_1+n_2)}$, for all $(x_1,x_2)\in M_1\times M_2$,
		
		\item[\emph{(d)}] $\operatorname{Scal}_M-\dfrac{\operatorname{Scal}_{M_2}}{f^2}\geq \operatorname{Scal}_{M_1}$, 
		
		\item[\emph{(e)}] $\operatorname{Scal}_{M_1}\leq 0$,
		
		\item[\emph{(f)}] $\operatorname{Scal}_{M_1}\geq\dfrac{n_1\operatorname{Scal}_{M_2}(x_2)}{n_2f^2(x_1)}$. 
  %and $n_1\geq n_2$.
	\end{itemize} 
\end{theorem}	

\begin{proof}
	{(a)} For $(x_1, x_2) \in M$ availing \eqref{f10},
	\begin{equation*}
		-f(x_1)\Delta f(x_1) + (n_2-1)\|\nabla f(x_1)\|^2+ \lambda(x_1,x_2)f^2(x_1)=\mu (x_2). 
	\end{equation*}
Then for $x_2=\Tilde{x_2}$, the above equation implies,
     \begin{equation*}
     	-f(x_1)\Delta f(x_1) + \lambda(x_1,\Tilde{x_2})f^2(x_1)=\mu (\Tilde{x_2})-(n_2-1)\|\nabla f(x_1)\|^2\leq 0,
     \end{equation*} 
 by using the hypothesis (a) which can be written as $\|\nabla f(x_1)\|\geq \sqrt{\frac{\mu(\Tilde{x_2})}{n_2-1}}$. This yields,
 \begin{equation*}
 	f(x_1)\Delta f(x_1)\geq \lambda(x_1,\Tilde{x_2})f^2(x_1).
 \end{equation*}
 Hence $\Delta f>0$ so warping function $f$ is constant, as base manifold is compact. \\
 
 
 
 { (b)} By \eqref{1.4},
 
 	$$\operatorname{Scal}_{M_1}= n_1\lambda(x_1,x_2)-n_2\frac{\Delta f(x_1)}{f(x_1)},$$
  equivalently
 	$$ \operatorname{Scal}_{M_1}+n_2\lambda(x_1,x_2)=(n_1+n_2)\lambda(x_1,x_2)-n_2\frac{\Delta f(x_1)}{f(x_1)}.$$
Consequently,
 	$$\operatorname{Scal}_{M}-\operatorname{Scal}_{M_1}=n_2\left(\lambda(x_1,x_2)+\frac{\Delta f(x_1)}{f(x_1)}\right),$$
 
which implies
$$\dfrac{\Delta f(x_1)}{f}+ \lambda(x_1,x_2)\leq 0,$$ 
by using assumption (b).
Hence $\Delta f \leq 0$, so warping function $f$ is constant.\newline
 
 {(c)} By \eqref{1.5},
 \begin{align*}
 	\operatorname{Scal}_{M_2}=n_2\mu(x_2) \geq n_2 \lambda(x_1,x_2),
 \end{align*}
by availing the part of the supposition $(c)$ which is $\mu(x_2)\geq\lambda(x_1,x_2)$ for all $(x_1,x_2)\in M_1\times M_2$.
 
Now by using \eqref{1.4} and the above equation,

	$$\operatorname{Scal}_{M_1}+ \operatorname{Scal}_{M_2}\geq (n_1+n_2)\lambda(x_1,x_2)-n_2 \frac{\Delta f(x_1)}{f(x_1)},$$
 thus,
	$$\operatorname{Scal}_{M_1}+ \operatorname{Scal}_{M_2}\geq \operatorname{Scal}_M-n_2 \frac{\Delta f(x_1)}{f(x_1)},$$
 which gives,   
 $$n_2 \frac{\Delta f(x_1)}{f(x_1)}	\geq \operatorname{Scal}_M -(\operatorname{Scal}_{M_1}+\operatorname{Scal}_{M_2})\geq 0,$$
by hypothesis in $(c)$. 
Hence, $\Delta f\geq0$ so warping function $f$ is constant.\newline
 
 {(d)} By \eqref{1.6}, we have
 
 	$$\operatorname{Scal}_{M_1}+\dfrac{\operatorname{Scal}_{M_2}}{f(x_1)^2}=\operatorname{Scal}_M+ n_2(n_2 -1)\dfrac{\|\nabla f\|^2}{f(x_1)^2}-2n_2\frac{\Delta f(x_1)}{f(x_1)}.$$
  
  Hence,
  
 	$$\operatorname{Scal}_M-\left(\operatorname{Scal}_{M_1}+\dfrac{\operatorname{Scal}_{M_2}}{f(x_1)^2}\right)=2n_2\frac{\Delta f(x_1)}{f(x_1)}-n_2(n_2 -1)\dfrac{\|\nabla f\|^2}{f(x_1)^2}\geq 0,$$
  
 which employing hypothesis in (d), 
  
 	$$2n_2\frac{\Delta f(x_1)}{f(x_1)}\geq n_2(n_2 -1)\dfrac{\|\nabla f\|^2}{f(x_1)^2}\geq 0.$$ 

This yields, $\Delta f\geq 0$ so warping function $f$ is constant.\newline

{(e)} By \eqref{1.4}, 
\begin{align*}
	n_2\dfrac{\Delta f(x_1)}{f(x_1)}&= n_1\lambda(x_1,x_2)-\operatorname{Scal}_{M_1}\geq 0,
\end{align*}
by hypothesis in $(e)$.
That is, $\Delta f\geq 0$ so warping function $f$ is constant.\newline

{(f)}
By \eqref{1.4} we have,

		$$\operatorname{Scal}_{M_1} = n_1\lambda(x_1,x_2)- n_2\dfrac{\Delta f(x_1)}{f(x_1)},$$

  thus,
		$$\operatorname{Scal}_{M_1} f^2(x_1) = n_1\lambda(x_1,x_2)f^2(x_1)- n_2f(x_1)\Delta f(x_1).\\$$

Now by using \eqref{f10},
$$\operatorname{Scal}_{M_1}f^2(x_1)-n_1\mu(x_2)=(n_1-n_2)f(x_1)\Delta f(x_1)-n_1(n_2-1)\|\nabla f(x_1)\|^2\geq 0,$$
by hypothesis in $(f)$.
Finally, we affirm that 
$$(n_1-n_2)f(x_1)\Delta f(x_1) \geq n_1(n_2-1)\|\nabla f(x_1)\|^2\geq 0.$$

Here the following three cases occur.\\
{ (i)} $n_1 = n_2$, then clearly 
 $\|\nabla f(x_1)\| = 0,$  and the warping function $f$ must be constant
 as base manifold is connected. \newline
{ (ii)} $n_1 \geq n_2$, then
clearly, $\Delta f \geq  0,$ so warping function $f$ is constant.\newline
{ (iii)} $n_1 \leq n_2$, then 
 $\Delta f \leq 0 $ consequently, warping function $f$ is constant.
\end{proof}

\end{document}
